\underline{\textbf{Первая глава}} представляет собой обзор литературы по теме диссертационной работы. В начале главы приведены основные этапы развития нелинейной оптики. Далее рассмотрены н-о процессы третьего порядка, в особенности одномодовое четырёхволновое смешение (ОМЧВС) и межмодовое ЧВС (ММЧВС), а также кратко изложена теория поперечных мод оптических волокон (ОВ). Затем рассмотрены нелинейно-оптические процессы второго порядка в ОВ на примере ГВГ. Представлена история экспериментальных наблюдений и попыток теоретического объяснения данного эффекта. Приведены укороченные уравнения ГВГ и физические модели этого явления, включая фотогальваническую модель и модель самоорганизации экситонов с переносом заряда. Рассмотрены работы, посвящённые измерению КФС при ГВГ в ОВ. В заключительной части обзора описаны методы лазерной калориметрии, ПРЛК, а также экспоненциальная, импульсная, градиентная методики обработки температурных кинетик, используемые при определении коэффициента оптического поглощения кристаллов.